\documentclass[border=10pt]{standalone}
\usepackage{tikz}
\usetikzlibrary{math}
\usepackage{tikz-3dplot}
\begin{document}
\makeatletter
\tikzset{%
  view/.style 2 args={%  observer longitude and latitude (y upwards)
                      %  Remark. lomg=0 means x=0
    z={({-sin(#1)}, {-cos(#1)*sin(#2)})},
    x={({cos(#1)}, {-sin(#1)*sin(#2)})},
    y={(0, {cos(#2)})},
    evaluate={%
      \tox={sin(#1)*cos(#2)};
      \toy={sin(#2)};
      \toz={cos(#1)*cos(#2)};
    },
    longitude = #1,
    latitude = #2
  },
  sun/.style n args={3}{% longitude, latitude, light contrast in [0, 1]
    sun longitude = #1,
    sun latitude = #2,
    contrast = #3,
    evaluate={%
      real \sunx, \suny, \sunz, \lightC;
      \sunx = sin(\sLongit)*cos(\sLatit);
      \suny = sin(\sLatit);
      \sunz = cos(\sLongit)*cos(\sLatit);
    }
  }
}
\pgfkeys{/tikz/.cd,
  latitude/.store in=\aLatit,  % observer's latitude
  latitude=0
}
\pgfkeys{/tikz/.cd,
  longitude/.store in=\aLongit,  % observer's longitude
  longitude=0  % corresponds to x=0
}
\pgfkeys{/tikz/.cd,
  sun latitude/.store in=\sLatit,
  sun latitude = 80
}
\pgfkeys{/tikz/.cd,
  sun longitude/.store in=\sLongit, 
  sun longitude=90
}
\pgfkeys{/tikz/.cd,
  contrast/.store in=\lightC,  % light contrast \in [0, 1]
  contrast=.5
}
\tikzmath{%
  function hornFaceColor(\i, \j, \N, \M) {%
    % Verifies if seen when the cone has no base; the color number is
    % modified by -2000.
    real \ang;
    \p@x = 1/\rB +(\i-.5)/\N*(1/\rT -1/\rB);
    \r@ = 1/\p@x;
    \t = 360*((\j-.5)/\M);
    \ux = \vx*cos(\t) +\wx*sin(\t);
    \uy = \vy*cos(\t) +\wy*sin(\t);
    \uz = \vz*cos(\t) +\wz*sin(\t);
    % modification needed when the radii are different
    \ang = atan(\r@*\r@); % angle with (Oy)
    \ux = \ux*cos(\ang) +\nx*sin(\ang);
    \uy = \uy*cos(\ang) +\ny*sin(\ang);
    \uz = \uz*cos(\ang) +\nz*sin(\ang);    
    \res = \ux*\tox + \uy*\toy + \uz*\toz;
    \tmp = int(100*\lightC*(\ux*\sunx + \uy*\suny + \uz*\sunz));
    if \res>0 then {% if seen
      return \tmp;
    } else {return {-2000 +\tmp};};    
  };
}
\makeatother
\tikzset{%
  pics/horn/.style args={xB=#1, xT=#2}{%
    % The interval for which the horn will be drawn (the curve y=1/x
    % rotated around the (Ox) axis).
    code={%
      \colorlet{mainColor}{.}
      \colorlet{innerColor}{-mainColor!50!gray}
      \colorlet{leftRGB{1}}{white}
      \colorlet{leftRGB{0}}{mainColor!50!black}
      \colorlet{leftRGB{3}}{white}
      \colorlet{leftRGB{2}}{innerColor!50!black}
      \tikzmath{%  
        real \rB, \rT;
        \rB = 1/#1;
        \rT = 1/#2;
        integer \Nx, \Nyz, \k, \j, \prevj, \i, \previ;
        \Nx = 50;
        \Nyz = 64;
        real \tmpBT, \tmpx, \tmpy, \tmpz, \tmpcst, \cstForColor;
        real \px, \nx, \ny, \nz, \vx, \vy, \vz, \wx, \wy, \wz;
        \nx = 1;
        \ny = 0;
        \nz = 0;
        \vx = 0;
        \vy = 1;
        \vz = 0;
        \wx = 0;
        \wy = 0;
        \wz = 1;
        %% points \P{\i,\j}
        % \i for x direction and
        % \j for directions in the plane (Oyz)
        for \i in {0, ..., \Nx}{%
          \px = #1 +\i/\Nx*(#2 -#1);
          \r = 1/\px;
          for \j in {0, ..., \Nyz}{%
            \t = \j/\Nyz*360;
            \Px{\i,\j} = \px +\r*\vx*cos(\t) +\r*\wx*sin(\t);
            \Py{\i,\j} = \r*\vy*cos(\t) +\r*\wy*sin(\t);
            \Pz{\i,\j} = \r*\vz*cos(\t) +\r*\wz*sin(\t);
          };
        };
        %% faces
        for \i in {1, ..., \Nx}{%
          \previ = \i -1;
          for \j in {1, ..., \Nyz}{%
            \prevj = \j -1;
            \cstForColor = hornFaceColor(\i, \j, \Nx, \Nyz);
            if \cstForColor>-999 then {%
              if \cstForColor>=0. then { \k = 1; } else {%
                \k = 0;  % in the shade for leftRGB{}
                \cstForColor = int(abs(\cstForColor));
              };
              {%
                \filldraw[leftRGB{\k}!\cstForColor!mainColor, line width=.001pt]
                (\Px{\previ,\prevj}, \Py{\previ,\prevj}, \Pz{\previ,\prevj})
                -- (\Px{\previ,\j}, \Py{\previ,\j}, \Pz{\previ,\j})
                -- (\Px{\i,\j}, \Py{\i,\j}, \Pz{\i,\j})
                -- (\Px{\i,\prevj}, \Py{\i,\prevj}, \Pz{\i,\prevj})
                -- cycle;
              };            
            } else {%
              \cstForColor = -(\cstForColor +2000);  % back to -value
              if \cstForColor>=0. then {%
                \k = 3;
                \cstForColor = \cstForColor;
              } else {%
                \k = 2;  % in the shade for leftRGB{}
                \cstForColor = int(abs(\cstForColor));
              };
              {%
                \begin{pgfonlayer}{background} 
                  \filldraw[leftRGB{\k}!\cstForColor!innerColor,
                  line width=.001pt]
                  (\Px{\previ,\prevj}, \Py{\previ,\prevj}, \Pz{\previ,\prevj})
                  -- (\Px{\previ,\j}, \Py{\previ,\j}, \Pz{\previ,\j})
                  -- (\Px{\i,\j}, \Py{\i,\j}, \Pz{\i,\j})
                  -- (\Px{\i,\prevj}, \Py{\i,\prevj}, \Pz{\i,\prevj})
                  -- cycle;
                \end{pgfonlayer}
              };
            };
          };
        };
      }  % end tikzmath
    }
  },
}


\begin{tikzpicture}[view={-30}{15}, sun={-50}{45}{.75}]
  \pgfdeclarelayer{background}
  \pgfsetlayers{background,main}
  
  \draw[->] (0, 0, 0) -- (2, 0, 0) node[shift={(.2, 0, 0)}, above] {$x$};
  \draw[->] (0, 0, 0) -- (0, 2, 0) node[shift={(0, .2, 0)}] {$y$};
  \draw[->] (0, 0, 0) -- (0, 0, 2) node[shift={(0, 0, .2)}, left] {$z$};

  \path (0, 0) pic[black!80!red!70, scale=1.5]
  {horn={xB=.95, xT=7}};
\end{tikzpicture}
\end{document}